%%%%%%%%%%%%%%%%%%%%%%%%%%%%%%%%%%%%%%%%%
% !TEX program = pdflatex
% Medium Length Professional CV
% LaTeX Template
% Version 2.0 (8/5/13)
%
% This template has been downloaded from:
% http://www.LaTeXTemplates.com
%
% Original author:
% Trey Hunner (http://www.treyhunner.com/)
%
% Important note:
% This template requires the resume.cls file to be in the same directory as the
% .tex file. The resume.cls file provides the resume style used for structuring the
% document.
%
%%%%%%%%%%%%%%%%%%%%%%%%%%%%%%%%%%%%%%%%%
%
% Revised for graduate school applications.
% Search for "TODO" to find places where you should fill in real facts/numbers.
%
%%%%%%%%%%%%%%%%%%%%%%%%%%%%%%%%%%%%%%%%%

%----------------------------------------------------------------------------------------
%	PACKAGES AND OTHER DOCUMENT CONFIGURATIONS 
%----------------------------------------------------------------------------------------
 
\documentclass{resume} % Use the custom resume.cls style 
\usepackage[dvipsnames]{xcolor}
\usepackage{hyperref}
\usepackage[left=0.4 in,top=0.3 in,right=0.4 in,bottom=0.3in]{geometry} % Document margins
\newcommand{\tab}[1]{\hspace{.2667\textwidth}\rlap{#1}}
\newcommand{\itab}[1]{\hspace{0em}\rlap{#1}}
\name{Qingyang Zhou} % Your name 
\address{Tsinghua University, Haidian, Beijing} % Your address
%\address{123 Pleasant Lane \\ City, State 12345} % Your secondary addess (optional) 
\address{+86 13269109251 \\ qy-zhou22@mails.tsinghua.edu.cn} % Your phone number and email
% TODO: your old PDF did not show the ORCID line, so resume.cls may only print 3 address lines.
% If ORCID is missing in the compiled PDF, merge it into the line below.
\address{https://orcid.org/0009-0004-4448-7663}
\address{https://esperarr.github.io/CV/ \\ https://github.com/Esperarr}

\definecolor{TsinghuaPurple}{cmyk}{0.58,0.90,0,0}
\hypersetup{
  colorlinks=true,
  linkcolor=TsinghuaPurple,
  urlcolor=TsinghuaPurple,
  citecolor=TsinghuaPurple,
  pdftitle={Qingyang Zhou - Curriculum Vitae},
  pdfauthor={Qingyang Zhou}
}
\renewenvironment{rSection}[1]{
\sectionskip
\textcolor{TsinghuaPurple}{\MakeUppercase{#1}}
\sectionlineskip
\hrule
\begin{list}{}{
%\setlength{\leftmargin}{1.5em}
\setlength{\leftmargin}{0em}
\setlength{\topsep}{4pt}
\setlength{\parsep}{0pt}
}
\item[]
}{
\end{list}
}

\renewenvironment{rSubsection}[4]{
 {\bf #1} \hfill {#2}
 \ifthenelse{\equal{#3}{}}{}{
  \\
  {\em #3}
 }
 \begin{list}{$\cdot$}{
  \setlength{\leftmargin}{1.0em}
  \setlength{\itemsep}{0.15em}
  \setlength{\topsep}{0.15em}
  \setlength{\parsep}{0pt}
  \setlength{\partopsep}{0pt}
 }
}{
 \end{list}
 \vspace{0.2em}
}


\begin{document}  
\setlength{\parindent}{0pt}
\setlength{\parskip}{0pt}
\def\sectionskip{\vspace{8pt}}
\def\sectionlineskip{\vspace{1pt}}

%----------------------------------------------------------------------------------------
%	EDUCATION SECTION
%----------------------------------------------------------------------------------------

\begin{rSection}{Education}

{\bf Tsinghua University} \hfill
{September 2022 -- Expected 2027}\\
\href{https://www.xyc.tsinghua.edu.cn/index.htm}{Xinya College} \& \href{https://www.cs.tsinghua.edu.cn/}{Department of Computer Science and Technology}
\\
GPA: 3.6/4.0 % TODO: if you have a major GPA, last-two-years GPA, or rank (e.g. top X\%), add it here
\\
Selected coursework: Complex Analysis (A+), Media Computing (A+), Database Special Topic Training (A+), Introduction to Artificial Intelligence, Introduction to Machine Learning, Fundamentals of Computer Graphics. % TODO: add grades for the other courses, or keep as is

\end{rSection} 

%--------------------------------------------------------------------------------------
%   MANUSCRIPTS
%--------------------------------------------------------------------------------------
\begin{rSection}{Manuscripts} \itemsep -3pt
% Both papers were submitted but not accepted, so they are listed as unpublished manuscripts
% (do NOT write "under review", "submitted to publication" or "accepted").
{\textbf{``Curriculum Multi-Reward Reinforcement Learning for Customized Text-to-Image Generation''}\\
 \textit{Yuwei Zhou, Xin Wang, Hong Chen, Yipeng Zhang, \textbf{Qingyang Zhou}, Wenwu Zhu}\\
 \textit{Manuscript}\quad \href{https://openreview.net/pdf?id=kiGUqX6Gct}{OpenReview}\quad \href{https://esperarr.github.io/CV/assets/papers/Curriculum_Multi_Reward.pdf}{Paper PDF}}

 {\textbf{``Towards Mutually Illuminative Collaborative Human--AI Deliberate Discussion''}\\
 % TODO: "Zhenhui" is missing a surname; please complete it.
 \textit{Zipeng Zhang, Shiwei Wu, Ran Ran, \textbf{Qingyang Zhou}, Ranrui Ma, Jiaye Leng, Jian Zeng, Zhenhui, Chun Yu}\\
 \textit{Manuscript}\quad \href{https://esperarr.github.io/CV/assets/papers/Towards_Mutually_Illuminative_Collaborative_Human\%E2\%80\%93AI_Deliberate_Discussion.pdf}{Paper PDF}}
 
\end{rSection}

%----------------------------------------------------------------------------------------
%	RESEARCH EXPERIENCE (reverse chronological by end date)
%----------------------------------------------------------------------------------------

\begin{rSection}{Research Experience}

\begin{rSubsection}{Research Intern, Information Retrieval Lab}{October 2024 -- July 2025}{\href{http://www.thuir.cn/}{Tsinghua University Information Retrieval Lab}}{}
\item Surveyed recommendation-model literature and reproduced open-source models (primarily FM and DeepFM), studying their implementations.
\item Integrated an existing sparse autoencoder (SAE) plug-in into the codebase and ran preliminary experiments under a senior student's guidance.
\end{rSubsection}  

\begin{rSubsection}{Research Intern, Pervasive Human-Computer Interaction Lab}{January 2025 -- May 2025}{\href{https://pi.cs.tsinghua.edu.cn/}{Tsinghua University Pervasive Human-Computer Interaction Laboratory}}{}
\item Surveyed and presented literature on aligning human and AI thought processes; contributed to group discussions on the research direction.
\item Implemented part of the core logic for a human--AI interactive writing assistant that supports collaborative idea generation.
\item Co-authored ``Towards Mutually Illuminative Collaborative Human--AI Deliberate Discussion'' (see Manuscripts).
\end{rSubsection}  

\begin{rSubsection}{Research Intern, Multimedia and Network Big Data Lab}{July 2023 -- July 2024}{\href{https://mn.cs.tsinghua.edu.cn/index/index.html}{Tsinghua University Multimedia and Network Big Data Lab}}{}
\item For the project ``Video Generation System Combining LLMs with Diffusion Models,'' reproduced experiments from multimodal LLM papers and presented results at group meetings.
\item Assisted senior students in debugging experimental code; studied diffusion models and personalized image generation.
\item Co-authored ``Curriculum Multi-Reward Reinforcement Learning for Customized Text-to-Image Generation'' (see Manuscripts).
\end{rSubsection}  

\end{rSection}

%----------------------------------------------------------------------------------------
%	INDUSTRY EXPERIENCE
%----------------------------------------------------------------------------------------

\begin{rSection}{Industry Experience}
 
\begin{rSubsection}{Client Development Engineer Intern, Tencent (Project MAX)}{May 2026 -- August 2026}{Tencent, Project MAX Team}{}
\item Optimized NavMesh-based character movement in battle, resolving stuttering on terrain bumps and sticking at convex corners; fixed a defect that allowed characters to leave the map during dash displacement by restricting destinations to the navigable region.
\item Developed a Scene-view debugging tool for the fixed-point logic layer that reads movement state at runtime through reflection and visualizes target points, input direction, and trajectories, with automatic stuck detection.
\item Developed in-battle GM tools, including performance overlays for the 2V2 mode and a buff inspector for the test scene; implemented the captain-selection hero showcase with model fallback, animation blending, and camera control; documented the NavMesh baking workflow.
\end{rSubsection}  
 
\end{rSection}

%-------------------------------------------------------------------------------
%	PROJECTS (reverse chronological)
%-------------------------------------------------------------------------------

\begin{rSection}{Selected Projects}

\begin{rSubsection}{Path Tracing Renderer \href{https://github.com/Esperarr/TinyPathTracer}{\small[Source code]}}{June 2026}{Computer Graphics Course Project}{}
\item Developed a physically based renderer supporting Whitted-style ray tracing and Monte Carlo path tracing with Russian roulette termination.
\item Implemented physically based materials with GGX microfacet importance sampling (isotropic and anisotropic), the Cook--Torrance BRDF, and Schlick's Fresnel approximation, covering diffuse, mirror, glass, and glossy surfaces.
\item Added BVH acceleration for meshes with over 100k faces, next-event estimation (NEE) with multiple importance sampling (MIS, power heuristic), depth of field, MSAA, gamma correction, and OpenMP parallelization.
\end{rSubsection}

%------------------------------------------------

\begin{rSubsection}{Five-Stage Pipeline CPU \href{https://github.com/Esperarr/thinpad-pipelined-core}{\small[Source code]}}{November 2025 -- December 2025}{Computer Organization Course Project}{}
\item Implemented a five-stage pipelined RISC-V CPU in SystemVerilog covering 19 basic instructions plus extensions, with data forwarding, control-signal generation, and hazard handling; validated through Vivado simulation.
\end{rSubsection}

%-------------------------------------------------

\begin{rSubsection}{Connect Four AI \href{https://github.com/Esperarr/Connect4_AI}{\small[Source code]}}{June 2024}{Course Project}{}
\item Developed a Connect Four AI in Python using Monte Carlo Tree Search (MCTS) with Upper Confidence bounds applied to Trees (UCT), achieving a win rate above 90\% against baseline AI opponents.
% TODO: specify the baseline opponents, number of games, and MCTS iteration budget.
\end{rSubsection}

%-------------------------------------------------

\begin{rSubsection}{Text Sentiment Analysis Model \href{https://github.com/Esperarr/SentimentAnalysis}{\small[Source code]}}{May 2024}{Course Project}{}
\item Implemented and compared CNN, RNN, and MLP text sentiment classifiers in PyTorch on a course-provided dataset, achieving over 90\% accuracy.
\end{rSubsection}

%-------------------------------------------------

\begin{rSubsection}{Intelligent Agent-Based Game ``Generals''}{September 2023 -- May 2024}{Major Project in the 28th Tsinghua University Intelligent Agent Competition}{}
\item Lead developer of a turn-based strategy game for the competition, available at \href{https://www.saiblo.net/game/35}{saiblo.net/game/35}.
\item Responsible for backend logic, unit testing, debugging, SDK development, and documentation; supported participants throughout the competition.
\end{rSubsection}  

\end{rSection}

%-----------------------------------------------------------------------------
%	TEACHING AND SERVICE
%-----------------------------------------------------------------------------

\begin{rSection}{Teaching and Service}

\begin{rSubsection}{Volunteer Teaching Assistant (Programming and Algorithms)}{September 2025 -- Present}{\href{https://learning.tsinghua.edu.cn/zxgk/ywjs.htm}{Tsinghua University, Center for Student Learning and Development}}{}
\item Provide drop-in tutoring in C/C++, data structures, and algorithms, covering debugging, algorithm design, and core concepts.
\end{rSubsection}

%------------------------------------------------

\begin{rSubsection}{Secretary, Intelligent Agent Department}{February 2023 -- Present}{\href{https://net9.org/home/}{Tsinghua Student Science and Technology Association}}{}
\item Participate in game development and event organization for the Tsinghua Intelligent Agent Competition; led development of ``Generals'' for the 28th edition (see Projects).
\end{rSubsection}  

%------------------------------------------------

\begin{rSubsection}{Tutor, Tsinghua University Science and Engineering Summer Camp}{July 2024 -- August 2024}{2024 Science and Engineering Summer Camp, I6 Class Tutor}{}
\item Provided technical guidance and mentored teams in intelligent robot design and cardboard-bridge construction.
\end{rSubsection}

\end{rSection}

%----------------------------------------------------------------------------------------
%	Skills
%----------------------------------------------------------------------------------------

\begin{rSection}{Skills}

% TODO: if you are not comfortable being asked about Java, JavaScript/TypeScript or Rust in an interview, remove them.
\begin{tabular}{ @{} >{\bfseries}l @{\hspace{3ex}} >{\raggedright\arraybackslash}p{0.68\textwidth} }
Languages & Chinese (native), English (TOEFL 108), Spanish (basic)\\
Programming & Python, C/C++, C\#, SystemVerilog, Java, JavaScript/TypeScript, Rust\\
AI/Machine Learning & PyTorch, NumPy, fundamentals of machine learning and deep learning\\
Graphics/Game Development & OpenGL, Unity, path tracing, rendering pipelines\\
Tools & Git, Linux, Docker, LaTeX, Markdown\\
 
\end{tabular}   

\end{rSection}

\end{document}